\section{Related work}
Network-on-Chip has been widely studied as a communication fabric for manycore, heterogeneous, and chiplet-based systems. In heterogeneous systems, different types of processing elements, such as CPUs, GPUs, memories, and accelerators, may have different communication requirements. For example, CPU-related traffic is often latency-sensitive, while GPU-related traffic usually requires higher throughput. This makes the design of the on-chip communication network more challenging than in a homogeneous system. Hybrid NoC architectures combining wireline and wireless links have been explored to improve communication efficiency in CPU-GPU heterogeneous manycore systems, especially for workloads with intensive data movement such as CNN training \cite{8119941}.

NoC design has also been investigated in the context of chiplet-based manycore architectures. Flexible NoC architectures with adaptable routers and links have been proposed to dynamically form different sub-network topologies, such as mesh, concentrated mesh, torus, and tree \cite{9218654}. These topologies can be selected according to application traffic patterns and different latency or bandwidth requirements. Such studies show that adapting the network structure can improve system performance and energy efficiency. However, this flexibility also introduces additional hardware and control complexity.